\section{Observation and KPI Spaces}

\begin{definition}[Enterprise State Space]
Let \(\Omega := \Omega(K_{EA})\) denote the OC/ETS enterprise state space of the
enterprise continuum \(K_{EA}\).
All constructions in this section are carried out strictly within the Ontology
of Continua and the frozen executable specification
\texttt{ETS.K\_EA.v1.1}.
No additional primitives, regularity assumptions, or modeling hypotheses are
introduced.
\end{definition}

\begin{definition}[Observable Domain \(\Omega_{\mathrm{obs}}\)]
The executable specification \texttt{ETS.K\_EA.v1.1} includes
(i) an information-axis component \emph{observability},
(ii) a factor \(S_{\mathrm{obs}}\) in the factorization of the continuum measure
\(k_{EA}\), and
(iii) an explicit \emph{ObservabilityInvariant}.
Accordingly, define the observable domain
\[
\Omega_{\mathrm{obs}}
\;:=\;
\bigl\{
\omega \in \Omega
\;\big|\;
S_{\mathrm{obs}}(\omega)\ \text{is defined and the ObservabilityInvariant holds}
\bigr\}.
\]
Only states in \(\Omega_{\mathrm{obs}}\) are admissible inputs to observation.
States for which the ETS phase condition \texttt{observability\_lost} applies lie
outside the scope of this definition.
\end{definition}

\begin{definition}[Observation Space \(\Sigma_{\mathrm{obs}}\)]
Define the observation space \(\Sigma_{\mathrm{obs}}\) as the codomain of all
OC/ETS-admissible readout maps defined on \(\Omega_{\mathrm{obs}}\).
Once a finite KPI family \(K=(\kappa_1,\ldots,\kappa_m)\) is fixed,
\(\Sigma_{\mathrm{obs}}\) is identified with a subset of \(\mathbb{R}^m\).
\end{definition}

\begin{definition}[KPI Family and Readout Map]\label{def:kpi-family}
Fix a finite KPI family
\(K := (\kappa_1,\ldots,\kappa_m)\) such that each
\(\kappa_i:\Omega_{\mathrm{obs}}\to\mathbb{R}\) is an OC/ETS-admissible readout,
i.e.\ it is defined on \(\Omega_{\mathrm{obs}}\) and does not violate the
ObservabilityInvariant.
The induced KPI readout map is
\[
\pi_{K}:\Omega_{\mathrm{obs}}\to\Sigma_{\mathrm{obs}}\subseteq\mathbb{R}^m,
\qquad
\pi_{K}(\omega)
=
\bigl(\kappa_1(\omega),\ldots,\kappa_m(\omega)\bigr).
\]
No assumption of injectivity, completeness, faithfulness, sufficiency, or
statistical optimality is imposed.
\end{definition}

\begin{definition}[\(K\)-Indistinguishability and Fibers]
Two states \(\omega_1,\omega_2\in\Omega_{\mathrm{obs}}\) are
\emph{\(K\)-indistinguishable} if and only if
\[
\pi_{K}(\omega_1)=\pi_{K}(\omega_2).
\]
Equivalently, they lie in the same fiber of \(\pi_K\), defined as
\[
\mathrm{Fib}_K(y)
\;:=\;
\pi_K^{-1}(y)
\;\subseteq\;
\Omega_{\mathrm{obs}},
\qquad
y\in\Sigma_{\mathrm{obs}}.
\]
This fiber representation is standard in observation and identification theory
\cite{kalman1960,ljung1999}.
\end{definition}

\begin{definition}[Identifiability on a Domain]
Let \(D\subseteq\Omega_{\mathrm{obs}}\) be any domain of interest.
KPI-based observation is said to be \emph{identifying on \(D\)} if the restricted
map \(\pi_K|_{D}\) is injective, i.e.,
\[
\forall \omega_1,\omega_2\in D:\quad
\pi_K(\omega_1)=\pi_K(\omega_2)
\;\Longrightarrow\;
\omega_1=\omega_2.
\]
If \(\pi_K|_{D}\) fails to be injective, then there exist distinct states in \(D\)
that are KPI-indistinguishable.
This definition corresponds to structural identifiability in the classical
system identification sense \cite{ljung1999,walter1997}.
\end{definition}

\begin{definition}[Quotient Induced by KPI Observation]
The equivalence relation \(\sim_K\) induced by KPI equality defines the quotient
space
\[
\Omega_{\mathrm{obs}} \big/ \!\sim_K ,
\]
whose elements are the KPI fibers.
KPI-based observation factors through this quotient: all structural variation
within a fiber is necessarily lost once observation is restricted to KPI
readouts.
This loss of information is intrinsic to non-injective observation maps
\cite{ash1990}.
\end{definition}
